\section{Neuromorphic Hardware and Energy Modelling}
\label{sec:background_neuromorphic}

\subsection{Neuromorphic Hardware Platforms}

Intel Loihi 2 \cite{gralewicz2024spikes} and IBM TrueNorth \cite{merolla2014million} represent the two principal reference platforms for neuromorphic computing.
Unlike conventional \glspl{gpu}, which execute computation synchronously across all neurons at every clock cycle, both architectures operate asynchronously and event-driven.
A neuron consumes energy only when it receives a spike and crosses its firing threshold, remaining entirely silent, drawing no dynamic power in the absence of input activity.
This mirrors the \gls{snn}'s spike-based dynamics precisely, making neuromorphic hardware the natural target deployment environment for event-driven architectures and completing the end-to-end alignment with the event camera described in Section~\ref{subsec:bg_snn_event_pairing}.
The per-\gls{sop} energy costs of these platforms, $16.1$\,pJ for Intel Loihi 2 \cite{gralewicz2024spikes} and $26.0$\,pJ for IBM TrueNorth \cite{merolla2014million}, form the hardware extrapolation targets used in the energy evaluation of this project.

\subsection{The MAC vs AC Energy Distinction}
\label{subsec:bg_mac_ac}

The foundational energy analysis underlying this project rests on the silicon-level operation costs documented by \textcite{horowitz20141}.
A standard \gls{cnn} performs a full \gls{mac} at each synaptic connection - a floating-point multiplication followed by an accumulation, costing $E_{\text{MAC}} = 4.6$\,pJ per 32-bit operation.
An \gls{snn}, by contrast, communicates exclusively via binary spikes; since a spike is either $0$ or $1$, the multiplication is eliminated entirely, reducing each synaptic event to an \gls{ac} operation at a cost of $E_{\text{AC}} = 0.9$\,pJ - a $5\times$ reduction in per-operation energy at the silicon level \cite{horowitz20141}.
This asymmetry is the first of two complementary mechanisms underpinning the \gls{snn}'s energy advantage over its \gls{cnn} counterpart; the second, input-dependent activation sparsity, is introduced and materialised alongside the complete per-frame energy model in Sections~\ref{subsec:dynamic_power} and~\ref{subsec:energy_consumption}.